\begin{thebibliography}{99}

\bibitem[Aas et al.(2021)]{aas}
Aas, K., Jullum, M. and L{\o}land, A. (2021). Explaining individual predictions when features are
dependent: More accurate approximations to Shapley values. \emph{Artificial Intelligence}, 298, 103502.
\href{https://doi.org/10.1016/j.artint.2021.103502}{doi:10.1016/j.artint.2021.103502}.

\bibitem[Athey et al.(2019)]{grf}
Athey, S., Tibshirani, J. and Wager, S. (2019). Generalized random forests.
\emph{Annals of Statistics}, 47(2), 1148--1178.
\href{https://doi.org/10.1214/18-AOS1709}{doi:10.1214/18-AOS1709}.

\bibitem[Breiman(2001)]{breiman}
Breiman, L. (2001). Random forests. \emph{Machine Learning}, 45, 5--32.
\href{https://doi.org/10.1023/A:1010933404324}{doi:10.1023/A:1010933404324}.

\bibitem[CFPB(2026a)]{regb}
Consumer Financial Protection Bureau. Regulation B, 12 CFR \S1002.9, Notifications.
Accessed September 22, 2026.
\url{https://www.consumerfinance.gov/rules-policy/regulations/1002/9/}.

\bibitem[CFPB(2026b)]{interp}
Consumer Financial Protection Bureau. Official Interpretations to Regulation B,
Comments 9(b)(2)-1 through -9. Accessed September 22, 2026.
\url{https://www.consumerfinance.gov/rules-policy/regulations/1002/interp-9/}.

\bibitem[CFPB(2025)]{withdrawn}
Consumer Financial Protection Bureau. Withdrawn Guidance. Entries for Circulars
2022-03 and 2023-03, withdrawn May 12, 2025. Accessed September 22, 2026.
\url{https://www.consumerfinance.gov/compliance/guidance/withdrawn-guidance/}.

\bibitem[Chen and Guestrin(2016)]{xgboost}
Chen, T. and Guestrin, C. (2016). XGBoost: A scalable tree boosting system.
\emph{Proceedings of ACM SIGKDD}, 785--794.
\href{https://doi.org/10.1145/2939672.2939785}{doi:10.1145/2939672.2939785}.

\bibitem[Chen et al.(2020)]{truedata}
Chen, H., Janizek, J. D., Lundberg, S. and Lee, S.-I. (2020). True to the model or true to the data?
\emph{arXiv preprint}, arXiv:2006.16234.
\url{https://arxiv.org/abs/2006.16234}.

\bibitem[Covert et al.(2021)]{covert}
Covert, I., Lundberg, S. and Lee, S.-I. (2021). Explaining by removing: A unified framework for
model explanation. \emph{Journal of Machine Learning Research}, 22(209), 1--90.

\bibitem[Good(1950)]{good}
Good, I. J. (1950). \emph{Probability and the Weighing of Evidence}. London: Charles Griffin.

\bibitem[Janzing et al.(2020)]{janzing}
Janzing, D., Minorics, L. and Bl{\"o}baum, P. (2020). Feature relevance quantification in
explainable AI: A causal problem. \emph{Proceedings of Machine Learning Research}, 108, 2907--2916.

\bibitem[Ke et al.(2017)]{lightgbm}
Ke, G., Meng, Q., Finley, T., Wang, T., Chen, W., Ma, W., Ye, Q. and Liu, T.-Y. (2017).
LightGBM: A highly efficient gradient boosting decision tree.
\emph{Advances in Neural Information Processing Systems}, 30.

\bibitem[Lei et al.(2018)]{loco}
Lei, J., G'Sell, M., Rinaldo, A., Tibshirani, R. J. and Wasserman, L. (2018).
Distribution-free predictive inference for regression.
\emph{Journal of the American Statistical Association}, 113(523), 1094--1111.
\href{https://doi.org/10.1080/01621459.2017.1307116}{doi:10.1080/01621459.2017.1307116}.

\bibitem[Lundberg and Lee(2017)]{lundberg}
Lundberg, S. M. and Lee, S.-I. (2017). A unified approach to interpreting model predictions.
\emph{Advances in Neural Information Processing Systems}, 30.

\bibitem[Lundberg et al.(2020)]{treeshap}
Lundberg, S. M., Erion, G., Chen, H., DeGrave, A., Prutkin, J. M., Nair, B.,
Katz, R., Himmelfarb, J., Bansal, N. and Lee, S.-I. (2020).
From local explanations to global understanding with explainable AI for trees.
\emph{Nature Machine Intelligence}, 2, 56--67.
\href{https://doi.org/10.1038/s42256-019-0138-9}{doi:10.1038/s42256-019-0138-9}.

\bibitem[Nair et al.(2022)]{nair}
Nair, V. N., Feng, T., Hu, L., Zhang, Z., Chen, J. and Sudjianto, A. (2022).
Explaining adverse actions in credit decisions using Shapley decomposition.
\emph{arXiv preprint}, arXiv:2204.12365.
\url{https://arxiv.org/abs/2204.12365}.

\bibitem[OpenML(2022)]{heloc}
OpenML. Dataset 45023, \texttt{heloc}, version 1: the FICO Home Equity Line of Credit data
released for the FICO Explainable Machine Learning Challenge.
\url{https://www.openml.org/d/45023}. Accessed October 14, 2026.

\bibitem[Platt(1999)]{platt}
Platt, J. (1999). Probabilistic outputs for support vector machines and comparisons to
regularized likelihood methods. In \emph{Advances in Large Margin Classifiers}, 61--74. MIT Press.

\bibitem[Prokhorenkova et al.(2018)]{catboost}
Prokhorenkova, L., Gusev, G., Vorobev, A., Dorogush, A. V. and Gulin, A. (2018).
CatBoost: Unbiased boosting with categorical features.
\emph{Advances in Neural Information Processing Systems}, 31.

\bibitem[Siddiqi(2006)]{siddiqi}
Siddiqi, N. (2006). \emph{Credit Risk Scorecards: Developing and Implementing Intelligent Credit
Scoring}. Hoboken, NJ: Wiley.

\bibitem[Sudjianto(2026a)]{aarcgbdt}
Sudjianto, A. (2026a). Evidence-based adverse action reasons for gradient-boosted credit models:
Comparable passing applicants and the effect of the reference distribution. Working manuscript,
September 22, 2026.

\bibitem[Sudjianto(2026b)]{aarcneural}
Sudjianto, A. (2026b). Evidence-based adverse action reasons for neural credit models with monthly
histories: Interpretable temporal features, passing references and exact attribution checks.
Working manuscript, September 2026.

\bibitem[Sundararajan and Najmi(2020)]{baseline}
Sundararajan, M. and Najmi, A. (2020). The many Shapley values for model explanation.
\emph{Proceedings of Machine Learning Research}, 119, 9269--9278.
\url{https://proceedings.mlr.press/v119/sundararajan20b.html}.

\end{thebibliography}
